% Pregunta 4. Inflación y competitividad precio (bloque III.1-III.2) · peso alto
\section*{Pregunta 4. Inflación y competitividad precio}

\begin{enunciado}
\textbf{4.} Analice la evolución de la \textbf{inflación} desde 2021 y sus principales determinantes. ¿Qué factores podrían explicar que en 2026 la inflación española se sitúe por encima de la de la zona del euro? ¿Aprecia riesgo de efectos de segunda vuelta? Valore la evolución de la \textbf{competitividad precio y coste} a lo largo de la serie.

\smallskip
\textbf{Peso alto} · \tiempo{21} · guía: III.1 (índices, subyacente, efecto base, origen, segunda ronda), III.2 (tipo de cambio real), II.1 (CLU y margen), B.1 (un brote de inflación) · cuadros 3, 5 y 6; gráficos 6, 7 y 8.
\end{enunciado}

\begin{examen}
\textit{España ha sufrido 2 choques de precios de origen energético: el de 2021-2023, del que salió antes que la zona del euro, y el de 2026, en el que su inflación (5,0~\% en septiembre) supera a la europea (3,8~\%) sobre todo porque su inflación subyacente es más alta. Hay riesgo de efectos de segunda ronda, y la competitividad en costes ganada en 2008-2018 se ha erosionado en parte.}

\rotulo{a) La inflación desde 2021 y sus determinantes}
El brote de 2021-2023 fue, en su origen, un \textbf{choque de oferta importado}: la energía sube un 21,1~\% en 2021 y un 27,6~\% en 2022 y el deflactor de las importaciones, un 22,4~\% en 2022 (cuadro~5). Que el IPCA (8,3~\%) creciera mucho más que el deflactor del PIB (4,7~\%) confirma que la inflación venía de fuera, porque el IPC incluye lo importado y el deflactor no. El choque se trasladó después a los alimentos (+10,7 y +11,1~\% en 2022-2023) y a la subyacente (3,8 y 4,1~\%), los \emph{efectos indirectos}. Hubo también un componente de demanda: la reapertura y el ahorro embalsado en 2021-2022, con la brecha ya positiva desde 2022. Lo que \emph{no} hubo fue una espiral de precios y salarios: en 2022 la remuneración por asalariado creció un 4,7~\%, casi 4 puntos menos que los precios, de modo que los trabajadores absorbieron buena parte del choque (el salario real cayó un 3,6~\%), mientras los márgenes empresariales se ampliaban (el deflactor del PIB creció 2,6 puntos más que los CLU). Como importador neto de energía, España se empobreció frente al exterior: la relación real de intercambio empeoró en torno a un 6~\% en 2022.

España llegó antes al pico (10,7~\% en julio de 2022, frente al 10,6~\% de la zona del euro en octubre, gráfico~6) y bajó también antes: un 3,4~\% de media en 2023, frente al 5,4~\%, gracias a la caída de la energía ($-16{,}1$~\%). Lo explicarían un traslado más rápido de los precios mayoristas de la electricidad a los consumidores, en los 2 sentidos, y medidas como el tope al precio del gas. El \emph{efecto base} se ve con claridad en marzo de 2023: la inflación española cae del 6,0 al 3,1~\% en un mes porque sale de la comparación la subida de marzo de 2022. La subyacente, en cambio, tocó techo más tarde (5,1~\% a principios de 2023) y bajó despacio, por los servicios: la «última milla». En 2024-2025 la inflación española (2,9 y 2,7~\%) quedó algo por encima de la europea (2,4 y 2,1~\%), con una subyacente media del 2,8 y el 2,6~\%: la economía española llegó al choque de 2026 sin haber completado del todo la desinflación.

\rotulo{b) Por qué en 2026 la inflación española supera a la europea}
El nuevo choque es otra vez energético: el Brent pasa de 62,5 dólares en diciembre de 2025 a 117,3 en abril y 114,2 en septiembre de 2026 (gráfico~7), una subida de más del 80~\%, algo mayor en euros. La inflación española sube del 2,4~\% en enero al 5,0~\% en septiembre, y la europea, del 1,7 al 3,8~\%. Lo revelador es \textbf{dónde está la diferencia}: de los 1,2 puntos de diferencia en septiembre, 1,0 están en la subyacente (3,5~\% frente a 2,5~\%), mientras la distancia entre la inflación general y la subyacente, que aproxima el peso de la energía y los alimentos, es parecida (1,5 y 1,3 puntos). La diferencia, por tanto, \uline{no viene tanto del petróleo como de factores internos}:
\begin{itemize}
\item una \textbf{demanda más fuerte}: brecha positiva del 2~\% y un consumo que sigue tirando en 2026 (pregunta~1);
\item \textbf{costes laborales} que crecen más: CLU del 4,4~\% en 2025, frente al 3,3~\% en la zona (cuadro~3);
\item el peso de los \textbf{servicios} (turismo, hostelería) y de la \textbf{vivienda}, cuyo precio sube a 2 dígitos (pregunta~3);
\item y, en menor medida, un traslado más rápido de la energía a la factura eléctrica.
\end{itemize}
La previsión del 3,0~\% para 2026 ha quedado desfasada: la media de enero a septiembre ya es del 3,6~\%.

\rotulo{c) ¿Riesgo de efectos de segunda vuelta?}
\textbf{Sí, y mayor que en la zona del euro.} Los indicios: la subyacente española sube del 3,0~\% en enero al 3,6-3,7~\% en mayo-julio y se mantiene en el 3,5~\%, mientras la europea sigue en el 2,2-2,6~\%; los salarios ya crecían el 4,4~\% en 2025 con la productividad estancada; y la brecha y la tensión del mercado de trabajo dan poder de negociación. Además, las empresas llegan al choque con los márgenes estrechados: el deflactor del PIB creció menos que los CLU en 2025 ($\gamma_{1+\mu}\approx3{,}0-4{,}4=-1{,}4$~puntos), así que tenderán a trasladar el encarecimiento de la energía para recuperarlos. Si los salarios se indexan a la inflación pasada, se arma la espiral. Lo atenúan unas expectativas que, a juzgar por la subyacente europea, siguen ancladas, y la subida de tipos del BCE (pregunta~7). En la terminología de la guía, la inflación europea tiene rasgos de \emph{transitoria} (concentrada en la energía, con la subyacente contenida), y la española, de \emph{persistente} (subyacente al alza, salarios acelerados y brecha positiva). Con una política monetaria común que no puede responder a España, la respuesta adecuada sería un \textbf{pacto de rentas} que reparta la pérdida de renta real del choque energético entre salarios y beneficios, una política fiscal que no añada demanda (pregunta~9) y medidas que aumenten la competencia en los mercados de bienes y servicios, para que el choque no se traslade con margen.

\rotulo{d) La competitividad precio y coste a lo largo de la serie}
La serie tiene 3 tiempos (cuadros~3 y~6, gráfico~8). Hasta 2008, \textbf{pérdida}: el tipo de cambio efectivo real con CLU frente a la zona del euro (un aumento es una apreciación) sube a 111,9. Entre 2008 y 2018, \textbf{devaluación interna}: sin tipo de cambio propio, la competitividad se recupera por los costes; los CLU crecen un 1,2~\% acumulado en 2007-2013, frente al 12,1~\% en la zona, y el tipo real cae casi un 16~\% hasta 94,4 en 2018. Desde 2019, \textbf{erosión parcial}: los CLU suben un 26,8~\% en 2019-2025, frente al 23,0~\% en la zona, y el tipo real se aprecia casi un 5~\% desde su mínimo de 2018. En el conjunto de 2007 a 2025 España conserva una ganancia de competitividad en costes (el tipo real con CLU cae un 9~\%), mientras que con precios de consumo frente a 42 socios apenas varía ($-2$~\%). La prueba de que la competitividad ha funcionado son las exportaciones en volumen, un 56~\% por encima de 2007 (gráfico~8). El riesgo es que el diferencial de inflación y de costes de 2025-2026 se prolongue: \uline{en una unión monetaria, cada punto de inflación por encima de los socios es competitividad que se pierde sin que nadie devalúe}.
\end{examen}

\begin{complemento}{4.1 · Por qué se contesta así: de dónde viene y si persiste}
\que{Qué se aplica.} Guía, III.1: IPC frente a deflactor (inflación importada); la subyacente «dice si un choque de oferta se está extendiendo»; el efecto base, $\pi^{12}_t-\pi^{12}_{t-1}\approx\pi^m_t-\pi^m_{t-12}$; el triángulo de \autor{Gordon} (demanda, costes, inercia) \parencite{Gordon2011Phillips}; directos, indirectos y segunda ronda; transitoria o persistente y qué política. III.2: la convención del tipo de cambio real. II.1: CLU y margen.

\begin{tablacaja}
\centering\scriptsize
\caption{El triángulo de \autor{Gordon} con los datos de España en 2025-2026}\label{tab:gordon}
\begin{tabular}{>{\raggedright\arraybackslash}p{2.6cm}>{\raggedright\arraybackslash}p{7.2cm}>{\raggedright\arraybackslash}p{4.2cm}}
\toprule
\rowcolor{tceeazulpalido}
\textbf{Origen} & \textbf{Indicio en el cuadernillo} & \textbf{Lectura} \\
\midrule
Costes (oferta) & Brent +83~\% en dólares y +86~\% en euros (dic. 2025-sep. 2026, G7); CLU +4,4~\% (C3) & El detonante; común con la zona en la energía \\
Demanda & Brecha +2,0~\% (C1); consumo privado +0,67~pp en el 2T de 2026 (G1); vivienda +12,7~\% (C4) & Propio de España: explica el diferencial \\
Inercia & Subyacente del 3,0 al 3,5-3,7~\% (G6); salarios +4,4~\% con la productividad estancada (C3) & El riesgo de persistencia \\
\bottomrule
\end{tabular}
\fuente{Cuadernillo (C: cuadro; G: gráfico). Clasificación: guía, III.1 \parencite{Gordon2011Phillips}.}
\end{tablacaja}

\que{¿Transitoria o persistente?} Con la regla de la guía (III.1): \emph{transitoria} si se concentra en la energía con la subyacente contenida y salarios moderados; \emph{persistente} si la subyacente sube, los salarios se aceleran y la brecha es positiva. España tiene rasgos de la segunda; la zona del euro, de la primera. Eso justifica que la respuesta sea más intensa en España que en la zona, y como la monetaria es común, por la vía fiscal y de rentas.

\que{Error típico.} Atribuir el diferencial de 2026 «al petróleo» sin mirar la subyacente; o decir que el BCE debe «subir más por España»: la política monetaria responde a la zona del euro.
\end{complemento}

\begin{complemento}{4.2 · Los cálculos: diferencial, margen y competitividad}
\que{Descomponer el diferencial de septiembre de 2026} (gráfico~6): general $5{,}0-3{,}8=1{,}2$; subyacente $3{,}5-2{,}5=1{,}0$; general menos subyacente, $5{,}0-3{,}5=1{,}5$ en España y $3{,}8-2{,}5=1{,}3$ en la zona. No son contribuciones exactas (la subyacente no pesa lo mismo en los 2 índices), pero bastan para ver que el diferencial es sobre todo subyacente. Medias de enero-septiembre: general 3,6 frente a 2,8; subyacente 3,4 frente a 2,4.

\que{¿Es alcanzable el 3,0~\% previsto para 2026?} La media anual es la de los 12 meses: hasta septiembre suman 32,5, y para que la media fuera 3,0 los 3 últimos meses tendrían que promediar $(36-32{,}5)/3\approx1{,}2~\%$, lejos del 5,0~\% de septiembre.

\que{Inflación importada en 2022}: IPCA $-$ deflactor del PIB $=8{,}3-4{,}7=3{,}6$~puntos; relación real de intercambio, deflactores de exportación e importación: $1{,}146/1{,}224-1\approx-6{,}4~\%$ (empeora) y en 2023 $1{,}015/0{,}955-1\approx+6{,}3~\%$ (se recupera).

\begin{figuracaja}
\caption{España: CLU nominales, deflactor del PIB y margen, 2019-2025}\label{fig:s04}
\grafico[0.76\linewidth]{s04-margen}
\fuente{Cuadros~3 y~5 del cuadernillo (Eurostat). Margen: $\gamma_{1+\mu}\approx\gamma_P-\gamma_{\mathrm{CLU}}$ (guía, II.1). Qué hay que ver: en 2021-2022 los márgenes se amplían (el deflactor sube más que los CLU: las empresas trasladan más que el coste laboral) y la participación de los beneficios sube del 37,8 al 40,3~\% del VAB (cuadro~4); en 2025 se estrechan ($-1{,}4$): no hay «inflación de beneficios» al llegar el choque de 2026.}
\end{figuracaja}

\begin{tablacaja}
\centering\scriptsize
\caption{Costes y precios acumulados: España frente a la zona del euro (\%)}\label{tab:clu-acum}
\begin{tabular}{>{\raggedright\arraybackslash}p{4.4cm}*{4}{>{\centering\arraybackslash}p{1.9cm}}}
\toprule
\rowcolor{tceeazulpalido}
 & \textbf{2007-2013} & \textbf{2013-2019} & \textbf{2019-2025} & \textbf{2007-2025} \\
\midrule
CLU nominales, España & 1,2 & 5,0 & 26,8 & 34,7 \\
CLU nominales, zona del euro & 12,1 & 6,5 & 23,0 & 46,8 \\
IPCA, España & 13,4 & 3,4 & 21,5 & 42,6 \\
IPCA, zona del euro & 12,4 & 5,4 & 22,9 & 45,6 \\
Tipo de cambio efectivo real, CLU (variación) & $-10{,}0$ & $-1{,}9$ & 2,5 & $-9{,}4$ \\
\bottomrule
\end{tabular}
\fuente{Cálculo propio con los cuadros~3, 5 y~6 (\texttt{calculos.txt}). Tipo de cambio: un descenso es una depreciación real (ganancia de competitividad). Umbral del procedimiento de desequilibrios: $\pm3~\%$ en 3 años en la zona del euro (guía, III.2); 2022-2025: $98{,}9/99{,}8-1=-0{,}9~\%$, sin alerta.}
\end{tablacaja}

\que{Alternativas de redacción.} «La desinflación española de 2023 fue más rápida porque su inflación era más energética; la de 2026 es más alta porque es más \emph{doméstica}». \que{Qué añadir con más tiempo.} La inflación desigual (los hogares de renta baja gastan más en energía y alimentos, guía III.1) y las ayudas fiscales: temporales, focalizadas y que mantengan el incentivo a ahorrar energía (guía, VI.2). \que{Qué recortar en el examen.} El efecto base, la relación real de intercambio y el detalle del tipo de cambio con IPC: bastan la subyacente, el CLU y los 3 tiempos de la competitividad.
\end{complemento}
